\section{产业链定位}
通鼎互联对应的不是“AI光模块单一链条”，而是“原材料与光棒工艺—光纤/光缆、通信电缆、电力电缆和通信设备—运营商、电网与项目客户”的综合线缆链条，并叠加安全系统和新能源业务。公司年报披露了光纤预制棒、光纤、光缆、通信电缆和电力电缆等一体化能力，这构成成本、交付和运营商渠道的基础，但并不自动等于高端AI订单。

\begin{exhibitbox}[表3-1\quad 全链条角色与估值边界]
\begin{tabularx}{\exhibitboxwidth}{L{2.6cm}L{2.7cm}X L{2.5cm}}
\textbf{链条节点} & \textbf{通鼎角色} & \textbf{已确认信息} & \textbf{估值边界} \\
上游材料/设备 & 成本输入 & 供应商和输入价格未逐项披露 & 不单独估值 \\
光纤/光缆 & 中游制造 & 2025收入、销量和库存有官方数据；H1称量价齐升 & 周期盈利信用，不给AI纯度溢价 \\
通信/电力电缆 & 规模底盘 & 收入和毛利率有官方披露 & 传统制造估值 \\
运营商/电网 & 需求锚点 & 集采和基础设施周期可做行业背景 & 不等于已确认订单 \\
安全系统 & 第二利润池 & 2025收入和毛利率披露，H1称经营稳健 & 需H1分部现金与利润验证 \\
\end{tabularx}
\sourcenote{analysis/full\_chain\_taxonomy.md；analysis/value\_chain\_economics.md；sources/official-20260714/}
\end{exhibitbox}

\section{竞争格局}
光纤光缆和通信电缆是规模、认证、原材料和运营商招标能力共同决定的行业。通鼎的优势在于光棒—光纤—光缆的链条完整和既有渠道；劣势在于2025年光纤光缆收入规模小于主流综合线缆平台，且公开披露不足以证明其在高端AI数据中心光纤中的客户份额、ASP或利用率。竞争优势可以支持周期修复估值，不能单独支持光模块龙头倍数。

\section{需求锚与证据边界}
中国移动特种光缆和普通光缆集采、运营商资本开支、电网投资和数据中心扩容构成行业需求背景。报告将这些内容标记为demand anchor；除非公司披露中标、合同金额、交付和收入确认，否则不把需求规模乘上假设份额写进通鼎利润预测。社交媒体中关于英伟达、亚马逊或特定AI订单的表述不进入证据链。
